\chapter{Bleeding During Pregnancy}

\section*{Definition}
Any vaginal bleeding occurring during pregnancy. Early light bleeding is common and is not always a sign of a problem, but bleeding should be reported to an obstetric clinician because causes range from benign cervical changes to pregnancy emergencies.

\section*{When to See a Doctor}

\begin{itemize}
 \item Contact your obstetric clinician about any bleeding. Seek emergency care for heavy bleeding, severe or one-sided abdominal pain, shoulder pain, faintness, fever, fluid leakage, regular contractions, or reduced fetal movement.
\end{itemize}

\section*{How It Develops}
First-trimester bleeding (most common) results from implantation bleeding, threatened miscarriage, or ectopic pregnancy. Second- and third-trimester bleeding raises concern for placenta previa (placenta covering the cervical os), placental abruption (premature separation of the placenta from the uterine wall), or vasa previa (fetal vessels crossing the membranes over the cervical os). Uterine rupture is a rare but catastrophic cause.

\section*{Causes}
\begin{itemize}
 \item Light early-pregnancy bleeding, including bleeding from a sensitive cervix; implantation bleeding is possible but cannot be diagnosed reliably from bleeding alone
 \item Threatened, inevitable, or complete miscarriage
 \item Ectopic pregnancy (fallopian tube)
 \item Molar pregnancy (gestational trophoblastic disease)
 \item Placenta previa
 \item Placental abruption
 \item Cervical insufficiency or incompetent cervix
 \item Cervical or vaginal infection (cervicitis, vaginitis)
 \item Vasa previa
 \item Uterine rupture (rare, usually scar-related)
\end{itemize}

\section*{Related Symptoms}
\begin{itemize}
 \item Abdominal pain
 \item Back pain
 \item Contractions
 \item Fever
 \item Decreased fetal movement
\end{itemize}


\section*{Medical Evaluation}
\begin{itemize}
 \item The clinician assesses vital signs, gestational age, bleeding amount and timing, pain, trauma, and fetal movement. Abdominal and speculum examinations are performed when appropriate; digital examination may be deferred until placenta previa is excluded later in pregnancy.
 \item Ultrasound is selected according to gestational age and symptoms to assess pregnancy location and viability, and later to assess placental location and fetal status.
 \item Tests may include a pregnancy hormone level in early pregnancy, complete blood count, blood type, Rh status, and infection testing. Rh immune globulin is considered according to gestational age, amount/cause of bleeding, Rh status, and local guidance; it is not automatically required for every Rh-negative patient.
 \item Referral to a maternal-fetal medicine specialist is arranged for complex cases such as placenta accreta, vasa previa, or recurrent bleeding of uncertain aetiology.
\end{itemize}

\section*{Treatment Options}
\begin{itemize}
 \item Management of threatened miscarriage depends on ultrasound and examination findings; routine bed rest has no proven benefit, and activity or intercourse advice is individualized.
 \item Ectopic pregnancy requires immediate medical or surgical intervention with methotrexate or salpingectomy depending on stability and size.
 \item Placenta previa and placental abruption require obstetric management, monitoring, and preparation for hemorrhage. Timing and route of delivery depend on bleeding, maternal and fetal condition, gestational age, placental findings, and local protocols.
\end{itemize}

\section*{Self-Care and Home Management}
\begin{itemize}
 \item Contact the maternity unit or obstetric clinician for advice about any bleeding; go to emergency or maternity triage immediately for heavy bleeding, severe pain, faintness, or reduced fetal movement.
 \item Do not insert anything into the vagina — no tampons, douching, or sexual intercourse until cleared by a healthcare provider.
 \item Use a pad to monitor the amount, colour, and presence of clots in the bleeding for reporting to the medical team.
 \item Sit or lie down while arranging care, use a pad to monitor bleeding, and call emergency services if bleeding is heavy, pain is severe, you feel faint, or fetal movement is reduced.
 \item Do not take any medications (including over-the-counter NSAIDs) without obstetric approval, as some are contraindicated during pregnancy.
\end{itemize}

\section*{References}
\begin{thebibliography}{99}
\bibitem{bleeding_during_pregnancy:ref1} Oyelese Y, Ananth CV. ``Placenta previa and placental abruption.'' \textit{N Engl J Med}. 2022;386(6):569-579.
\bibitem{bleeding_during_pregnancy:ref2} Hendriks E, MacNaughton H, et al. ``First-trimester bleeding: evaluation and management.'' \textit{Am Fam Physician}. 2019;99(3):166-174.
\bibitem{bleeding_during_pregnancy:ref3} Royal College of Obstetricians and Gynaecologists. ``Antepartum haemorrhage: green-top guideline No. 63.'' RCOG, 2011.
\bibitem{bleeding_during_pregnancy:ref4} Cunningham FG, Leveno KJ, et al. \textit{Williams Obstetrics}, 26th ed. McGraw-Hill, 2022.
\bibitem{bleeding_during_pregnancy:ref5} American College of Obstetricians and Gynecologists. ``Bleeding During Pregnancy.'' ACOG patient guidance, accessed September 2026. https://www.acog.org/womens-health/faqs/bleeding-during-pregnancy.
\bibitem{bleeding_during_pregnancy:ref6} Royal College of Obstetricians and Gynaecologists. ``Placenta Praevia and Placenta Accreta: Diagnosis and Management (Green-top Guideline No. 27a).'' 2025. https://www.rcog.org.uk/guidance/browse-all-guidance/green-top-guidelines/placenta-praevia-placenta-accreta-diagnosis-and-management-green-top-guideline-no-27a/.
\end{thebibliography}
